%==============================================================
% METHODOLOGY SECTION -- journal version.
%
% Rewritten against the conference version to add what the new results
% required: the image normalisation that removes the storage artefact, the
% article-level screen, the enlarged sealed partition, and the multi-seed
% threshold-free protocol.
%
% Labels defined here and referenced by the Results section:
%   sec:scope  sec:data  sec:normalise  sec:dedup  sec:articles
%   sec:partition  sec:features  sec:fusion(*)  sec:training
% (*) sec:fusion lives in the Results section; it is referenced from here.
%
% Verified facts, traceable to files in the project root:
%   split composition, article ids   BanFakeMM CSVs; doi:10.17632/3r9s7wfbh5.1
%   screening counts                 duplicate_clusters.csv, dedup_summary.txt
%   enlarged partition               big_sealed_results_v2.csv (n=1308)
%   article overlap                  article_screen.csv (86 of 1308)
%   seeds and statistics             adaptation_final_seeds.csv
%==============================================================

\section{Methodology}

\subsection{System Architecture and Scope}
\label{sec:scope}

\sysname{} scores an uploaded image from its visual and captioned content
alone. Three representations are extracted---a BLIP caption encoded by
multilingual SBERT and refined through Cross-Image Semantic Fusion (CISF), a
CLIP visual embedding, and TF-IDF features over the caption---and the
concatenated vector is classified by a gradient-boosted tree
(Fig.~\ref{fig:arch}). Section~\ref{sec:fusion} shows the caption pathway earns
nothing, so the configuration we recommend is the CLIP branch alone; the full
pipeline is described because it is what the ablation evaluates. A
retrieval-augmented verification stage exists in the implementation but is not
evaluated here.

One property of this design governs everything that follows: the system is
image-derived end to end. The BLIP caption, the OCR text and the CLIP embedding
are all read off the image, and the article's own title and body---which the
benchmark distributes---are never an input. The system is therefore well posed
only where the manipulation reaches the pixels. Section~\ref{sec:detectable}
establishes which of the benchmark's attack types satisfy that condition, and
we scope the evaluation to those.

\subsection{Benchmark and Attack Types}
\label{sec:data}

We use BanFakeMM \cite{banfakemm}, a detector-guided multimodal Bangla
benchmark built from 7{,}040 authentic news articles with images, drawn from
eight Bangla newspapers across eleven topical categories. The release provides
five adversarial sets---AdversarialEdit, AdversarialGen, LinguisticAttack,
RandomSwap and SemanticSwap---each containing 1{,}000 authentic and 1{,}000
manipulated samples, so every split is balanced by construction. Each record
carries an \texttt{article\_id}, a \texttt{source\_dataset} and a category, and
in the released master file 1{,}863 article identifiers carry both an authentic
image and its manipulated counterparts. Provenance is published, and we use it
rather than infer it (Section~\ref{sec:articles}).

Labels are deterministic: a sample is manipulated if the benchmark's
construction pipeline produced it, not if an annotator judged it so. That makes
the labels reproducible and the task well defined, but it also means a model
may learn to recognise the construction pipeline rather than the manipulation.
Section~\ref{sec:normalise} removes the part of that signature which lives
outside the image content; the part that survives resampling we cannot separate
from the manipulation itself, and we say so where it matters.

The classifier is trained on AdversarialEdit and evaluated under shift on
AdversarialGen. The three remaining splits manipulate article text and leave
the photograph untouched, which makes them undecidable for an image-only model
though presumably well posed for one that also reads the text;
Section~\ref{sec:detectable} demonstrates this rather than assuming it, and
reports what a model does on them.

\subsection{Image Normalisation}
\label{sec:normalise}

Before any feature is extracted, every image passes through one identical
re-encoding: resized to $384 \times 384$ with Lanczos resampling, converted to
RGB, saved as JPEG at quality 90, with no EXIF block carried over. Resolution,
aspect ratio, file format, colour mode and EXIF are therefore constant across
classes; image content is preserved.

This is not cosmetic. In AdversarialEdit all 1{,}000 manipulated images are
$1408 \times 768$ while the authentic images span many resolutions and formats,
and a classifier given only that metadata---no pixels at all---separates the
classes perfectly (Section~\ref{sec:artefact}). Without normalisation, any
reported accuracy is unidentifiable: it may be reading content, or it may be
reading the file header. We report the normalised figures as the primary
results and, where a comparison with the uncontrolled setting is informative,
give both and label them.

\subsection{Contamination Screening}
\label{sec:dedup}

\subsubsection{Image-level}
Filename comparison is not sufficient, since a manipulated image is an edited
copy of an authentic one and the two can carry unrelated names. Every image the
study uses is screened by content instead. Each is reduced to a 64-bit DCT
perceptual hash (pHash), a 64-bit difference hash (dHash) and its CLIP
embedding; a pair is a near-duplicate when either Hamming distance is at most
$2$ or the CLIP cosine similarity is at least $0.98$, and pairs are merged with
union-find. These thresholds sit in a clear gap: of 3.5 million candidate
pairs, 94 lie at pHash distance $0$ and only 102 at distance $6$ or below.
Duplicates are near-exact, and looser settings chain unrelated images into
implausible clusters.

Screening is applied in both directions that matter: within the training split,
to build cross-validation groups, and from every evaluation partition against
everything the models have seen---the base training set together with the
adaptation set.

\subsubsection{Article-level}
\label{sec:articles}
An image screen removes repeated pictures, not repeated stories: two different
photographs of the same event share entities, scene and caption vocabulary, and
different pixels do not rule that out. Because every image in this benchmark is
named by its source article, the check is exact rather than inferred---the
filename stem is the \texttt{article\_id} in each distributed set, which we
verified by content hash against the released master file.

We therefore report article-level overlap alongside the image-level screen
rather than folding one into the other, and we state the screening claim
accordingly: the pipeline removes image duplicates, and measures the residual
article overlap. Section~\ref{sec:articlescreen} gives the measured effect on
the conclusions.

\subsection{Partitioning}
\label{sec:partition}

AdversarialEdit supplies the training corpus: 2{,}000 images, which the
image-level screen resolves into 1{,}980 groups. Every fold in
Section~\ref{sec:ablation} is built with \texttt{StratifiedGroupKFold} over
those groups and asserted group-disjoint at run time.

AdversarialGen supplies the shift evaluation. Setting aside its distributed
80/10/10 partition, whose proportions do not suit a three-way adaptation study,
we repartition its 2{,}000 images at a fixed seed: 300 for supervised
adaptation, 150 for decision-threshold selection, 200 forming the sealed
partition used by earlier work on this system, and 1{,}350 held back entirely.
The adaptation and validation partitions are screened against the training set,
leaving 148 validation images; the adaptation partition keeps all 300, since a
duplicate inside training data is redundancy rather than leakage and pruning it
would confound the adaptation contrast with a change in training-set size.

The primary sealed partition is the block of 1{,}350 images, which entered no
training, tuning or threshold-selection step of any kind and had never been
scored. Screened against the base training set and the adaptation set under the
full criteria above---42 images removed, 21 by pHash, 40 by dHash and 24 by
CLIP cosine---it retains \textbf{1{,}308 images}, 658 of them manipulated. The
older 200-image block, which screening reduces to 194, has been scored before
and is therefore no longer pristine; we use it only where a direct comparison
with the earlier measurement is the point, and label it as such.

\subsection{Features and Classifier}
\label{sec:features}

A BLIP model generates a caption per image; OCR was evaluated during ablation
and removed (Section~\ref{sec:ablation}). The caption is tokenized and stemmed,
with separate handling for Bangla and Latin script, for TF-IDF features, and
separately encoded by multilingual SBERT. That embedding is refined through
CISF: combined with the weighted average of its $k=5$ nearest training
neighbours at fusion weight $\alpha=0.65$, re-normalized, and reduced to 32
principal components. A CLIP vision encoder
(\texttt{openai/clip-vit-base-patch32}) independently produces a
512-dimensional embedding, $L_2$-normalised and reduced to 64 components. These
were fixed design choices rather than outcomes of a search; a sensitivity study
over $k$ and $\alpha$, and comparison against adaptive fusion
\cite{mmfusionnet,claaf}, is left for future work.

The full vector concatenates four blocks, each fit only on the training fold:
the CISF caption embedding, the CLIP embedding, TF-IDF over the caption
($1$--$2$ grams, vocabulary capped at 5{,}000, minimum document frequency 3),
and auxiliary numeric features. Each training embedding's neighbour search
excludes itself, and at inference an unseen embedding is compared only against
this fixed training-set reference. Classification uses XGBoost
(\texttt{max\_depth}~3, \texttt{learning\_rate}~0.0742,
\texttt{n\_estimators}~443, \texttt{subsample}~0.866,
\texttt{colsample\_bytree}~0.683, \texttt{min\_child\_weight}~2,
\texttt{gamma}~0.143, \texttt{reg\_alpha}~$1.01\times10^{-5}$,
\texttt{reg\_lambda}~4.911), selected by randomized search over 8 sampled
configurations under 3-fold cross-validation.

\subsection{Protocol, Thresholds and Statistics}
\label{sec:training}

\subsubsection{Thresholds}
The operating threshold is selected on held-out validation data, never on an
evaluation partition, under a rule fixed before the sweep was run: $\tau$
maximizes fake-class F1 on the screened 148-image validation partition over a
$0.01$-step grid, ties broken by the median of the tied run. The sealed
partition is then scored once per threshold, after all training and threshold
decisions are frozen. Thresholds are selected separately for the normalised and
unnormalised settings, since the two produce differently calibrated scores, and
neither is reselected on test data.

Same-distribution figures are computed at the conventional $\tau = 0.50$
cut-off; no threshold search is performed on the training split, so those
numbers carry no selection effect.

\subsubsection{Cross-validation and dependence}
Ablation results come from repeated stratified 5-fold grouped cross-validation
over the clusters of Section~\ref{sec:dedup}---two repeats, ten folds---with
identical splits and fixed hyperparameters across every configuration and all
transforms fit inside the training fold only.

Repeated cross-validation does not yield independent observations: each image
is held out once per repeat, and any two training folds overlap in four fifths
of their rows. We therefore never stack repeats. Out-of-fold probabilities are
averaged across repeats so that each image contributes exactly one prediction
to any paired test, and fold-level comparisons use the Nadeau--Bengio corrected
resampled $t$ interval, which adds the train/test overlap ratio
($n_{\text{test}}/n_{\text{train}} = 0.25$ at five folds) to the naive $1/J$
variance term---a $3.5\times$ inflation here. The $\pm$ values reported beside
cross-validated accuracies are the spread across folds, not standard errors,
and no significance claim rests on them.

\subsubsection{Metrics and uncertainty}
Results are reported as accuracy, precision, recall and F1 for the manipulated
class, with authentic-class recall added under shift. Fake-recall is the more
safety-critical quantity---a missed manipulation circulates unchecked whereas a
false alarm costs only a review step---which is why the selection rule
maximizes fake-class F1.

For contrasts between two models on the same sealed partition we report AUROC
as the primary statistic. It is threshold-free, and the alternative---accuracy
at a fixed operating point---conflates a change in the model with a change in
where the cut-off happens to sit; Section~\ref{sec:adaptation} shows the two
disagreeing in sign on the same data. Confidence intervals on sealed-partition
quantities come from a paired bootstrap over images with $5{,}000$ resamples;
paired model comparisons additionally report an exact McNemar test over
discordant predictions.

\subsubsection{Seeds and reproducibility}
Every contrast on the sealed partition is run with eight seeds
($42$, $101$, $202$, $303$, $404$, $505$, $606$, $707$), varying fold
assignment, the PCA transforms and the classifier together; we report the
seed-averaged statistic and the seed-to-seed standard deviation, so that the
size of an effect can be read against the variability of the procedure that
produced it. Single-seed results are reported nowhere except where we compare
against an earlier single-seed measurement. All manifests, per-image
probabilities, screening clusters and the full validation sweeps are released
with the code.
